% ==============================================================================
%  Chapter 5 -- GELT: a gauge-equivariant lattice transformer
%  Purpose: THE methods chapter, ~20-25 pages. Complete, self-contained
%  definition of the architecture + equivariance proofs + design rationale.
%  Sources: notes/abstract.md, notes/GELTsummary.md, gelt/blocks_{rope,bias}.py,
%  gelt/lattice.py (transport). Proof details can go to Appendix B.
% ==============================================================================

\section{Design requirements}
\label{sec:gelt-requirements}

\TODO{State the contract up front: (i) exact gauge equivariance of every
hidden layer (adjoint transformation $W(x) \to \Omega(x) W(x) \Omega^{\dg}(x)$)
and exact invariance of the readout; (ii) configuration-dependent,
inspectable attention; (iii) receptive-field growth with loop content
sufficient for universality (the L-CNN argument must transfer);
(iv) group-generic code ($\mathbb{Z}_2 \to \SU{N}$ verbatim).}

\section{Parallel transport on the \texorpdfstring{$L^1$}{L1}-ball}
\label{sec:gelt-transport}

\TODO{The key geometric ingredient. Define the signed offset set
$1 \le |\Delta x|_1 \le R$; the transport $T_{\Delta x}(x)$ as the AVERAGE
over all shortest lattice paths (multinomially many), computed by a
$|\Delta x|_1$-ordered dynamic-programming recursion, not path
enumeration. Why averaging: restores lattice rotation symmetry that any
single canonical path breaks (the \texttt{mode="single"} variant exists
precisely to A/B this). State and prove gauge covariance
$T_{\Delta x}(x) \to \Omega(x) T_{\Delta x}(x) \Omega^{\dg}(x+\Delta x)$;
mention the octant identity
$T_{-\Delta x}(x) = T^{\dg}_{\Delta x}(x - \Delta x)$ as a consistency
check (unit-tested), deliberately NOT exploited at build time. Memory
cost table vs.\ $R$ and $D$ -- the wall that motivates offset-chunked
attention (\cref{sec:concl-future}).}

\section{The GEMHSA layer}
\label{sec:gelt-gemhsa}

\TODO{Walk the forward pass (mirrors notes/GELTsummary.md):
(1) augment the channel set with daggers and the identity;
(2) fused $Q, Q_v, K, V$ projections (channel-mixing only -- no color
mixing, so covariance is manifest);
(3) transport $K, V$ to the query site via $T$/$T^{\dg}$;
(4) gauge-INVARIANT score $s = \Re\Tr[Q^{\dg} \tilde K]$ + positional
encoding; (5) softmax over offsets;
(6) matrix-bilinear value $\sum_{\Delta x} \alpha \, Q_v^{\dg} \tilde V$
-- the L-Bilin analogue, which doubles loop length per layer and carries
the universality argument; (7) channel mix, residual, trace-gated
activation (L-Act). Figure: one block diagram.}

\subsection{Positional encodings: RoPE vs.\ offset bias}
\TODO{The two variants (\texttt{blocks\_rope} / \texttt{blocks\_bias},
$\approx 90\%$ shared code): rotary rotation of the score by $\Delta x$
vs.\ a learned per-offset bias. Both act only on the INVARIANT score, so
equivariance is untouched -- the geometric prior is decoupled from the
symmetry. Caveat to resolve before the final write-up: RoPE axis coverage
requires $d_{qkv} \ge 2D$ (small $d_{qkv}$ in 4D leaves axes unrotated).
\TODO{merge the two variants into one module with a
\texttt{pos\_encoding} switch and parametrize the equivariance tests over
both -- currently the TRAINED variant (rope) is the less-tested one.}}

\section{The full model}
\label{sec:gelt-model}

\TODO{\texttt{ChannelLift} $\to$ stacked GEMHSA blocks $\to$ \texttt{Trace}
$\to$ per-site MLP $\to$ spatial reduction (\texttt{mean} for global
targets, \texttt{"none"} for per-site/per-timeslice outputs -- the
glueball operator mode). Input: plaquette channels (optionally multiple
smearing levels, see \cref{sec:glue-run5}); parameter count vs.\ the
matched L-CNN and CNN baselines (table).}

\section{Equivariance verification}
\label{sec:gelt-verification}

\TODO{Three tiers, report all: (1) unit tests -- GEMHSA end-to-end
equivariance, $\SU{2}$ complex128 and $\mathbb{Z}_2$ float64 (bit-exact),
both gate branches, finite-grad backward on $\SU{3}$;
(2) full-model invariance check
$f(W^{\Omega}, T^{\Omega}) \approx f(W, T)$ on random $\SU{2}$ gauge
transformations; (3) \TODO{the worst-case-$\Omega$ stress test in float64
-- does not exist yet; run it and quote the drift}.}

\section{Computational cost}
\label{sec:gelt-cost}

\TODO{Scaling in $V$, $R$, heads, channels; attention over
$|L^1\text{-ball}|$ offsets per site; transport precompute vs.\
on-the-fly (the glueball training builds 3D transports per batch on the
fly). Honest table of wall-clock and memory on the V100 for the
production runs.}
